\documentclass[11pt,oneside]{article}
\usepackage[T1]{fontenc}
\usepackage[utf8]{inputenc}
\usepackage{amsmath,amsthm,mathtools}
\usepackage[dvipsnames]{xcolor}
\usepackage{newtxtext,newtxmath}
\usepackage[a4paper,margin=27mm,headheight=15pt]{geometry}
\usepackage{microtype,booktabs,array,longtable,enumitem}
\usepackage{fancyhdr}
\usepackage[most]{tcolorbox}
\usepackage{xurl}
\usepackage[colorlinks=true,linkcolor=MidnightBlue,citecolor=MidnightBlue,urlcolor=MidnightBlue]{hyperref}
\hypersetup{pdftitle={Hidden Negative Directions in Surcomplex Linear Algebra},pdfauthor={Research manuscript prepared for Vladimir Reshetnikov},pdfsubject={Hahn fields, algebraic approximation, restricted positivity, matrix-valued measures}}
\setlength{\parskip}{4pt}
\setlength{\parindent}{1.25em}
\setlist{nosep,leftmargin=2em}
\emergencystretch=2em
\pagestyle{fancy}
\fancyhf{}
\fancyhead[L]{\small Hidden negative directions}
\fancyhead[R]{\small Surreal and surcomplex research}
\fancyfoot[C]{\thepage}
\renewcommand{\headrulewidth}{0.3pt}
\newtheorem{theorem}{Theorem}[section]
\newtheorem{lemma}[theorem]{Lemma}
\newtheorem{proposition}[theorem]{Proposition}
\newtheorem{corollary}[theorem]{Corollary}
\theoremstyle{definition}
\newtheorem{definition}[theorem]{Definition}
\newtheorem{example}[theorem]{Example}
\newtheorem{question}[theorem]{Question}
\theoremstyle{remark}
\newtheorem{remark}[theorem]{Remark}
\newcommand{\R}{\mathbb R}
\newcommand{\C}{\mathbb C}
\newcommand{\Q}{\mathbb Q}
\newcommand{\Z}{\mathbb Z}
\newcommand{\N}{\mathbb N}
\newcommand{\No}{\mathbf{No}}
\newcommand{\K}{K_\Gamma}
\newcommand{\F}{F_\Gamma}
\newcommand{\Pfield}{\mathscr P_\Gamma}
\newcommand{\Efield}{\mathscr E_\Gamma}
\newcommand{\Cfield}{\mathscr C_\Gamma}
\newcommand{\supp}{\operatorname{supp}}
\newcommand{\rank}{\operatorname{rank}}
\newcommand{\diag}{\operatorname{diag}}
\newcommand{\Span}{\operatorname{span}}
\newcommand{\ord}{\operatorname{ord}}
\newcommand{\In}{\operatorname{In}}
\newcommand{\trdeg}{\operatorname{trdeg}}
\newcommand{\fg}{\mathrm{fg}}
\newcommand{\st}{\operatorname{st}}
\newcommand{\eps}{\varepsilon}
\newcommand{\primes}{\mathcal P}
\newcommand{\qprime}[1]{q_{#1}}
\newcommand{\restr}[2]{{#1}\!\upharpoonright_{#2}}
\newcommand{\repo}{\texttt{Surreal}}
\newcommand{\pin}{\texttt{048b72cf7cbfc8ab246e4f73788c10460cb3f6e0}}
\newtcolorbox{summarybox}{colback=MidnightBlue!3,colframe=MidnightBlue!45,boxrule=.5pt,arc=1mm,left=3mm,right=3mm,top=2mm,bottom=2mm,breakable}

\begin{document}
\begin{titlepage}
\centering
\vspace*{15mm}
{\small\sffamily RESEARCH MANUSCRIPT\par}
\vspace{9mm}
{\Huge\bfseries Hidden Negative Directions\par}
\vspace{3mm}
{\Huge\bfseries in Surcomplex Linear Algebra\par}
\vspace{9mm}
{\Large Prime-denominator approximation barriers\par}
\vspace{2mm}
{\Large and a two-scale matrix-measure criterion\par}
\vspace{13mm}
{\large Prepared for Vladimir Reshetnikov\par}
\vspace{2mm}
{22 September 2026\par}
\vspace{12mm}
\begin{minipage}{0.9\textwidth}
\begin{abstract}
Let $K_\Gamma=\C((t^\Gamma))$, with $\Gamma$ a nonzero divisible ordered
abelian group, and let $\mathscr P_\Gamma$ consist of those Hahn series
whose support lies in a finitely generated additive subgroup of $\Gamma$.
This is an algebraically closed subfield containing every monomial.
Nevertheless, for every $r\ge1$ we construct an explicit Hermitian matrix
of size $r+1$ which is strictly positive on every nonzero vector over
$\mathscr P_\Gamma$, and even over its valuation closure, but has inertia
$(1,r,0)$ over $K_\Gamma$.

The construction rests on prime-denominator tails. For
$h=\sum_{p\text{ prime}}t^{\delta(1-1/p)}$, every finite-lattice-supported
approximant differs from $h$ at valuation below $\delta$. Matching its
first $N$ terms by an element algebraic over $\C(t^\delta)$ requires,
and can be achieved with, degree exactly $\prod_{j=1}^N p_j$.
Almost-disjoint prime sets yield continuum many algebraically independent
tails and a uniform linear-combination separation estimate. These results
transport to actual surreal normal forms, including $t=\omega^{-1}$.

A second theorem characterizes positivity of a two-scale matrix-valued
measure $M_0+\eps M_1$ by ordinary scalar null ideals together with a
vector-valued absolute-continuity condition. The latter detects the
coupling into a residual kernel that scalar positivity misses.
All assertions proposed as contributions have written proofs. The
support-field construction, character automorphisms, finite spectral
algebra, and elementary density obstruction are explicitly treated as
background, not new discoveries.
\end{abstract}
\end{minipage}
\vfill
\begin{minipage}{.9\textwidth}
\small
\textbf{Status.} Candidate-original theorem package, not refereed or
formally verified. No named published conjecture is claimed settled.
The repository comparison is pinned and targeted, not exhaustive.
The accompanying finite checks are diagnostics, not proofs of the
infinite statements.
\end{minipage}
\end{titlepage}

\pagenumbering{roman}
\tableofcontents
\clearpage
\pagenumbering{arabic}

\section{The problem and the contribution}\label{sec:intro}

Positivity of a Hermitian matrix is a universal assertion about vectors:
\begin{equation}\label{eq:psd}
 A\succeq0\quad\Longleftrightarrow\quad z^*Az\ge0\quad\text{for every }z.
\end{equation}
Over the ordinary complex numbers, testing a dense collection of vectors
is sufficient, by continuity. Over a Hahn field, an algebraically large
collection need not be dense. Even an algebraically closed subfield
containing every monomial can miss an entire negative cone. This paper
makes that failure explicit and quantifies the obstruction.

The question is not whether finite-dimensional spectral theory survives
an ordered field extension. It does in the real-closed setting, and the
repository already develops that subject \cite{repo-spectral}. The
question is instead how much of a Hahn vector must be available before
quadratic-form testing detects the actual sign. Finite exponent groups,
algebraic extensions, and valuation limits are three natural sources of
probes. Our first theorem defeats all three simultaneously.

\subsection{A precise headline result}

We use increasing Hahn exponents, so $t^\gamma$ is infinitesimal when
$\gamma>0$. Write $\mathscr P_\Gamma$ for the union of the fields
$\C((t^H))$ over finitely generated additive subgroups $H\le\Gamma$.
The term \emph{finite-lattice support} always means membership in some
such $H$; it does not mean finite support or a discrete lattice.
For instance, $\Z+\sqrt2\Z$ is allowed when it is a subgroup of $\Gamma$.

\begin{theorem}[Maximal hidden negative inertia; preview]\label{thm:preview}
For each integer $r\ge1$, each nonzero divisible ordered abelian group
$\Gamma$, and each $\delta\in\Gamma_{>0}$, there are explicit real Hahn
series $h_1,\ldots,h_r$ of positive valuation such that, with
$h=(h_1,\ldots,h_r)^T$ and $\eps=t^\delta$,
\begin{equation}\label{eq:mainmatrix}
 A_r=\begin{pmatrix}
 1&h^T\\ h&hh^T-\eps^2 I_r
 \end{pmatrix}
\end{equation}
has the following properties.
\begin{enumerate}[label=(\roman*)]
\item $z^*A_rz>0$ for every nonzero
$z\in\mathscr P_\Gamma^{r+1}$, and also for every nonzero vector over the
valuation closure of $\mathscr P_\Gamma$ in $K_\Gamma$.
\item Over $K_\Gamma$, its inertia is $(1,r,0)$ and
$\det A_r=(-\eps^2)^r$.
\item The vectors $(-h_j,e_j)$ have quadratic value $-\eps^2$.
\item The $h_j$ can be chosen from one explicitly described family of
cardinality $2^{\aleph_0}$ that is algebraically independent over
$\mathscr P_\Gamma$.
\end{enumerate}
\end{theorem}

The proof appears in Sections~\ref{sec:tails}--\ref{sec:closure}.
The word \emph{hidden} is restricted to vector probes from the specified
subfield. It does not mean hidden from determinants, symbolic
elimination, or an algorithm allowed to use the entries of $A_r$ as new
scalars. Indeed, every principal $2\times2$ minor containing the first
coordinate is $-\eps^2$. The matrix is an extremely transparent
counterexample once its full entries are available.

\subsection{Why primes enter}

For the increasing sequence of primes $p_1,p_2,\ldots$, put
\begin{equation}\label{eq:alltail-intro}
 q_p=1-\frac1p,
 \qquad h=\sum_{p\in\primes}t^{q_p\delta},
 \qquad P_N=\prod_{j=1}^N p_j.
\end{equation}
A finitely generated exponent group containing $\delta$ meets
$\Q\delta$ in $\frac1m\Z\delta$ for some integer $m\ge1$.
A prime not dividing $m$ therefore supplies a coefficient of $h$ that
cannot be canceled by any series supported in that group.

More quantitatively, the least possible degree of an algebraic
approximant $a$ over $\C(t^\delta)$ satisfying
\[
 v(h-a)>q_{p_N}\delta
\]
is exactly $P_N$; see Theorem~\ref{thm:degree}. Thus the obstruction is
not merely that $h$ is transcendental. There is an exact field-degree
price for crossing each successive precision threshold.

\subsection{A matrix-valued measure result}

The scalar Hahn-measure null-ideal criterion is already present in the
repository. Its Herglotz report explicitly excludes a matrix-valued
null-ideal criterion from its scope \cite{repo-herglotz}. We resolve the
\emph{two-scale} version of that extension, not the full arbitrary-support
problem.

Let $M_0,M_1$ be ordinary finite Hermitian matrix-valued measures on a
measurable space. For an ordinary vector $u\in\C^d$, write
\[
 \mu_u=u^*M_0u,\qquad \nu_u=u^*M_1u,\qquad \sigma_u=M_1u.
\]
Theorem~\ref{thm:measure} proves that $M_0+\eps M_1$ is positive as a
Hahn-valued matrix measure exactly when $M_0$ is positive and, for every
ordinary $u$,
\begin{equation}\label{eq:measure-preview}
 \nu_u^-\ll\mu_u,
 \qquad
 |\sigma_u|\ll\mu_u+\nu_u^+.
\end{equation}
Here $\nu_u^\pm$ are the ordinary Jordan parts, $|\sigma_u|$ is ordinary
vector-measure total variation, and $\ll$ means absolute continuity of
null sets. The second condition is genuinely matrix-valued. It rules out
coupling into a direction whose leading and first residual quadratic
values both vanish.

\subsection{Novelty boundary and evidence}

The candidate contribution is the combined explicit package: the exact
primorial approximation law, the simultaneous separation estimate for
almost-disjoint prime tails over the entire finite-lattice field, maximal
hidden negative inertia persisting through valuation closure and workspace
enlargement, and the two-scale matrix-measure null-ideal formulation.
We did not locate these combined statements in the targeted repository
and literature review. That is not a certification of priority.

Hahn fields, support-defined subfields, character automorphisms, and
Puiseux constructions are established subjects
\cite{rayner,rayner-structures,automorphisms}. The necessary background is
proved or stated explicitly below. Nor is the bare assertion that a
nondense subfield can miss a negative quadratic value new mathematics:
Section~\ref{sec:density} isolates the elementary density principle so
that it is not confused with the quantitative construction.

The distinction between mathematics and evidence is important here.
The proofs, the exact finite checks, and any future Lean development are
three different things. This manuscript supplies the first two; it does
not supply a Lean verification. Appendix~\ref{app:audit} records the
repository snapshot and the limited scope of the searches.

\section{Hahn conventions and the probe field}\label{sec:hahn}

\subsection{Fields, valuation, conjugation, and positivity}

Unless stated otherwise, $\Gamma$ is a nonzero set-sized divisible
ordered abelian group. We set
\[
 F_\Gamma=\R((t^\Gamma)),\qquad K_\Gamma=\C((t^\Gamma)).
\]
An element is a formal sum $f=\sum_{\gamma\in S}f_\gamma t^\gamma$ with
well-ordered support $S\subseteq\Gamma$. For $f\ne0$,
$v(f)=\min S$; we set $v(0)=+\infty$. We use the usual conventions for
comparisons involving $+\infty$.

The order on $F_\Gamma$ is the sign of the leading real coefficient.
Complex conjugation acts coefficientwise and fixes precisely
$F_\Gamma$. For $z\in K_\Gamma$, write $|z|^2=z\bar z$; this notation
does not require choosing a square root. For $z\ne0$,
\begin{equation}\label{eq:normvaluation}
 |z|^2>0,\qquad v(|z|^2)=2v(z).
\end{equation}
If $z_1,\ldots,z_n$ are not all zero, positivity of the leading
coefficients gives
\begin{equation}\label{eq:sum-squares}
 v\!\left(\sum_{j=1}^n|z_j|^2\right)
 =2\min_j v(z_j).
\end{equation}
There is no cancellation among these leading coefficients.

The standard Hahn-field closedness theorem says that $K_\Gamma$ is
algebraically closed and $F_\Gamma$ is real closed under these
hypotheses. We use this established input, rather than claim another
proof of it; it belongs to the classical generalized-series theory
\cite{rayner}, and the repository separately develops the relevant Hahn
closedness machinery \cite{repo-map}.

A Hermitian matrix $A\in M_n(K_\Gamma)$ is positive semidefinite if
$z^*Az\ge0$ for all $z\in K_\Gamma^n$. Its inertia
$\In(A)=(n_+,n_-,n_0)$ records its numbers of positive, negative, and
zero squares under congruence over $K_\Gamma$. Finite-dimensional
Hermitian diagonalization over a real-closed base makes this
well-defined; the repository's spectral report treats this background
\cite{repo-spectral}.

\subsection{What summation means}

A family of Hahn series is \emph{strongly summable} when the union of
its supports is well ordered and each exponent receives a contribution
from only finitely many members. Products and the infinite sums used
here are understood in this sense. The usual positive-support lemma
licenses power-series substitution at positive-order elements
\cite{neumann}.

Strong summation is not asserted to be convergence of finite subsums in
the valuation topology. In particular, the exponents $q_p\delta$ in
\eqref{eq:alltail-intro} are all below $\delta$. Their terms form a valid
Hahn sum even though its finite partial sums do not approach it to
valuation greater than $\delta$. This bounded sequence of exponents is
precisely what the construction needs.

All field operations and supports are set-sized. When an actual surreal
application is made, each finite collection of inputs is placed in a
set-sized exponent workspace before any Hahn argument is applied.

\subsection{The finite-lattice-supported subfield}

\begin{definition}\label{def:probefield}
The complex and real finite-lattice-supported fields are
\begin{equation}\label{eq:probefields}
 \Pfield=\bigcup_{\substack{H\le\Gamma\\H\text{ finitely generated}}}
               \C((t^H)),
 \qquad
 \Efield=\bigcup_{\substack{H\le\Gamma\\H\text{ finitely generated}}}
               \R((t^H)).
\end{equation}
A \emph{probe vector} is a vector with coordinates in $\Pfield$.
\end{definition}

A finitely generated subgroup of an ordered abelian group is a finite-rank
free abelian group. It need not be convex, discrete, or cofinal.
The union in \eqref{eq:probefields} is directed: the sum of finitely many
finitely generated subgroups is finitely generated. As each
$\C((t^H))$ is a subfield of $K_\Gamma$, $\Pfield$ is a field.
It contains every complex constant and every monomial $t^\gamma$.
It is closed under arbitrary restrictions of a support, and
\[
 \Pfield=\Efield[i].
\]
This is an instance of the established support-family approach to Hahn
subfields, not a new field construction \cite{rayner-structures}.

\begin{example}\label{ex:puiseux}
When $\Gamma=\Q$, every finitely generated subgroup is
$\frac1n\Z$ for a suitable integer $n\ge1$, after harmless enlargement.
Consequently
\[
 \mathscr P_\Q=\bigcup_{n\ge1}\C((t^{1/n})),
\]
the ordinary complex Puiseux-series field. Its elements can have infinite
support, but each individual element has bounded denominators.
For general $\Gamma$, infinitely many exponents in a fixed rank-two or
rank-ten subgroup are just as admissible.
\end{example}

\begin{proposition}[Closure under positive-order formal evaluation]
\label{prop:evaluation}
Let $a_1,\ldots,a_m\in\Pfield$ have positive valuation. Every formal
series $\Phi\in\C[[X_1,\ldots,X_m]]$ has a defined Hahn evaluation
$\Phi(a_1,\ldots,a_m)$, and this value belongs to $\Pfield$.
\end{proposition}
\begin{proof}
Choose one finitely generated subgroup $H$ containing all the input
supports. The finite union of their positive supports is well ordered.
The positive-support lemma gives strong summability of all monomial
contributions, including finiteness over all word lengths. Every resulting
exponent is a finite sum of input exponents and hence lies in $H$.
The evaluated series therefore belongs to $\C((t^H))$.
\end{proof}

Thus infinitesimal exponential and logarithm do not leave the probe
field. This assertion concerns formal evaluation at positive-order
arguments; it is not a claim that every version of surreal exponentiation
preserves this field.

\section{Algebraic support and character orbits}\label{sec:characters}

This section supplies an algebraic support bound with all details needed
later. The underlying diagonal character automorphisms are standard in
Hahn-field automorphism theory \cite{automorphisms}. Related finite-index
monomial character arguments are also already used by the repository
\cite{repo-spectral}. They are tools here, not novelty claims.

\subsection{Extending characters}

\begin{lemma}[Character extension and separation]\label{lem:characterextension}
Let $B$ be a subgroup of an abelian group $A$. Every homomorphism
$B\to\C^\times$ extends to a homomorphism $A\to\C^\times$.
Moreover, characters separate the elements of any abelian group.
\end{lemma}
\begin{proof}
The group $\C^\times$ is divisible: every element has an $n$th root
for each $n\ge1$. Order the extensions of the given character by inclusion
of their domains. The union along a chain is an extension, so a maximal
one exists by Zorn's lemma.

If its domain $B'$ omits $a\in A$, consider $B'+\Z a$.
If no positive multiple of $a$ lies in $B'$, assign any nonzero complex
value to $a$. Otherwise let $n$ be the least positive integer with
$na\in B'$ and choose $c\in\C^\times$ with
$c^n=\chi(na)$. The prescription
$\chi(b+ka)=\chi(b)c^k$ is well-defined and extends the character.
Both cases contradict maximality. Hence the maximal domain is $A$.

For separation, define a character nontrivial on a given nonzero element
on the cyclic group it generates, using a primitive root of unity when
its order is finite and, for example, the value $2$ when its order is
infinite. Extend it to the entire group.
\end{proof}

For a subgroup $G\le\Gamma$ and a character
$\chi:\Gamma/G\to\C^\times$, define
\begin{equation}\label{eq:characteraction}
 T_\chi\!\left(\sum a_\gamma t^\gamma\right)
 =\sum a_\gamma\chi(\gamma+G)t^\gamma.
\end{equation}
The support is unchanged. Addition is immediate and multiplication
follows from $\chi(\alpha+\beta)=\chi(\alpha)\chi(\beta)$ in each
finite coefficient convolution. Thus $T_\chi$ is a field automorphism
with inverse $T_{\chi^{-1}}$, and it fixes $\C((t^G))$ pointwise.

\subsection{The support-index bound}

\begin{lemma}[Algebraic degree bounds the support index]\label{lem:supportdegree}
Let $G\le\Gamma$ be finitely generated. Suppose $f\in K_\Gamma$ is
algebraic over a field $L$ satisfying
\[
 \C(t^G)\subseteq L\subseteq\C((t^G)),
\]
and let $D=[L(f):L]$. Put $B=G+\langle\supp f\rangle$.
Then
\begin{equation}\label{eq:supportdegree}
 [B:G]\le D,
 \qquad
 \supp f\subseteq\frac1{D!}G.
\end{equation}
The notation $\frac1{D!}G$ is inside the divisible group $\Gamma$.
No assertion that $[B:G]$ divides $D$ is made.
\end{lemma}
\begin{proof}
Every $T_\chi(f)$ is a root of the same minimal polynomial over $L$;
there are at most $D$ distinct such images. Fix
$\gamma\in\supp f$. If $\gamma+G$ has infinite order, characters of
its cyclic subgroup give arbitrarily many different values at that
class, and Lemma~\ref{lem:characterextension} extends all of them.
The coefficient at $\gamma$ then gives more than $D$ distinct images of
$f$, a contradiction. If the class has finite order $m>D$, its $m$
possible character values give the same contradiction. Hence its order
is at most $D$, so $D!\gamma\in G$.

It follows that $B/G$ is a subgroup of
$(\frac1{D!}G)/G$, a finite group because $G$ is finitely generated.
Write $e=[B:G]$. A finite abelian group of order $e$ has exactly $e$
complex characters, as is seen by decomposing it into finite cyclic
groups. Each character extends from $B/G$ to $\Gamma/G$.
Two different characters of $B/G$ differ on at least one support
exponent, because the classes of those exponents generate $B/G$.
Their images of $f$ are therefore different. Thus $e\le D$.
\end{proof}

\begin{remark}
The second bound in \eqref{eq:supportdegree} is deliberately coarse.
The first bound remembers the entire finite quotient generated by the
support. In the prime construction it is the first bound, not the
factorial estimate, that produces the sharp answer.
\end{remark}

\begin{proposition}[Algebraic closedness of the probe field]
\label{prop:probeclosed}
For divisible $\Gamma$, $\Pfield$ is algebraically closed and
$\Efield$ is real closed. In particular every element of $K_\Gamma$
algebraic over the monomial field
\[
 \C(t^\gamma:\gamma\in\Gamma)
\]
belongs to $\Pfield$.
\end{proposition}
\begin{proof}
A polynomial over $\Pfield$ has finitely many coefficients. Choose a
single finitely generated $G$ containing all their supports. A root $f$
exists in $K_\Gamma$ by the standard Hahn closedness theorem.
It is algebraic over $\C((t^G))$, so
Lemma~\ref{lem:supportdegree} places its support in
$\frac1{D!}G$, again a finitely generated group. Thus $f\in\Pfield$.
Every nonconstant polynomial over $\Pfield$ therefore has a root in it.

The real field $\Efield$ is ordered, and
$\Efield[i]=\Pfield$ is algebraically closed. The usual
characterization of real-closed ordered fields gives real closedness.
Alternatively, positive square roots can be chosen in $F_\Gamma$ and
then in $\Efield$ by the support bound, while an odd-degree polynomial
has a real root in $F_\Gamma$ to which the same bound applies.

Finally the monomial field is contained in $\Pfield$. A polynomial over
it splits in $\Pfield$, so every one of its roots in $K_\Gamma$ is
already in $\Pfield$.
\end{proof}

This proposition strengthens the interpretation of a probe. We are not
merely testing rational expressions or bounded-degree algebraic
expressions. We permit an algebraically closed field, containing all
monomial scales at once and arbitrary Hahn series in every individual
finitely generated exponent group.

\section{A prime-denominator approximation barrier}\label{sec:barrier}

Fix $\delta>0$ and set $s=t^\delta$. Let $\primes$ be the ordinary set
of positive primes, in increasing order. Define
\begin{equation}\label{eq:hdefinition}
 h_\delta=\sum_{p\in\primes}s^{1-1/p},\qquad
 h_{\delta,N}=\sum_{j=1}^N s^{1-1/p_j}.
\end{equation}
The support of $h_\delta$ is well ordered of order type $\omega$.
Its valuation is $\delta/2$. Every support exponent is strictly less
than $\delta$.

\subsection{A denominator lemma}

\begin{lemma}\label{lem:cyclicintersection}
If $H\le\Gamma$ is finitely generated and contains $\delta$, then
\begin{equation}\label{eq:cyclicintersection}
 H\cap\Q\delta=\frac1m\Z\delta
\end{equation}
for some integer $m\ge1$.
If $p$ is prime and $p\nmid m$, then $q_p\delta\notin H$, and for every
other prime $q$,
\begin{equation}\label{eq:freshcoset}
 q_p\delta+H\ne q_q\delta+H.
\end{equation}
\end{lemma}
\begin{proof}
The intersection is a subgroup of the finitely generated abelian group
$H$ and hence is finitely generated. As a subgroup of the rational line
$\Q\delta$, it is cyclic. Since it contains $\delta$, its positive
generator is $\delta/m$ for an integer $m\ge1$.

The rational number $1-1/p$ has denominator $p$ in lowest terms.
If $p\nmid m$, it does not belong to $\frac1m\Z$. For distinct primes
$p,q$, the difference
\[
 q_p-q_q=\frac{p-q}{pq}
\]
has a denominator divisible by $p$ after reduction, since $p\nmid q$
and $p\nmid(p-q)$. It therefore also fails to belong to
$\frac1m\Z$, proving \eqref{eq:freshcoset}.
\end{proof}

\begin{theorem}[Uniform finite-lattice approximation barrier]
\label{thm:barrier}
For every $a\in\Pfield$,
\begin{equation}\label{eq:barrier}
 v(h_\delta-a)<\delta.
\end{equation}
In particular $h_\delta\notin\Pfield$, and no element algebraic over
the field of all monomials can approximate $h_\delta$ to valuation
$\delta$ or higher.
\end{theorem}
\begin{proof}
Choose a finitely generated group $H$ containing $\delta$ and
$\supp a$. Let $m$ be as in Lemma~\ref{lem:cyclicintersection} and choose
a prime $p\nmid m$. The coefficient of $a$ at $q_p\delta$ is zero,
whereas that of $h_\delta$ is one. Thus $h_\delta-a$ is nonzero and
\[
 v(h_\delta-a)\le q_p\delta<\delta.
\]
The last assertion follows from Proposition~\ref{prop:probeclosed}.
\end{proof}

\subsection{The exact field-degree cost}

\begin{theorem}[Primorial approximation law]\label{thm:degree}
For $N\ge1$, let $P_N=\prod_{j=1}^N p_j$. Among all $a\in K_\Gamma$
algebraic over $\C(s)$ and satisfying
\begin{equation}\label{eq:precision}
 v(h_\delta-a)>q_{p_N}\delta,
\end{equation}
the minimum possible degree $[\C(s,a):\C(s)]$ is exactly $P_N$.
The finite sum $h_{\delta,N}$ attains this minimum and satisfies
\begin{equation}\label{eq:partialprecision}
 v(h_\delta-h_{\delta,N})=q_{p_{N+1}}\delta.
\end{equation}
The same minimum holds with $\R(s)$ and real approximants.
\end{theorem}
\begin{proof}
Let $a$ satisfy \eqref{eq:precision} and let
$D=[\C(s,a):\C(s)]$. At every exponent $q_{p_j}\delta$ for $j\le N$,
the coefficient of $a$ must be one. Apply
Lemma~\ref{lem:supportdegree} with $G=\Z\delta$.
The subgroup of $\Gamma/G$ generated by these support classes contains
\[
 \frac{\delta}{p_1}+G,\ldots,\frac{\delta}{p_N}+G,
\]
since $q_p\delta\equiv-\delta/p\pmod G$.
These classes generate $(\frac1{P_N}\Z\delta)/(\Z\delta)$, whose order
is $P_N$. Therefore $D\ge P_N$.

On the other hand, $h_{\delta,N}\in\C(s^{1/P_N})$, an extension of
$\C(s)$ of degree at most $P_N$. Its support contains the first $N$
exponents, so the same lower bound shows that its degree is exactly
$P_N$. Its difference from $h_\delta$ has least exponent
$q_{p_{N+1}}\delta$, which proves both admissibility and
\eqref{eq:partialprecision}.

For real approximants, complexifying can only decrease field degree,
so the same lower bound applies. The displayed real finite sum has
real degree at most $P_N$ and hence exactly $P_N$.
\end{proof}

\begin{center}
\begin{tabular}{rrrr}
\toprule
$N$ & $p_N$ & Precision crossed, divided by $\delta$ & Least degree $P_N$\\
\midrule
1&2&$1/2$&2\\
2&3&$2/3$&6\\
3&5&$4/5$&30\\
4&7&$6/7$&210\\
5&11&$10/11$&2310\\
6&13&$12/13$&30030\\
\bottomrule
\end{tabular}
\end{center}

The inequalities in the theorem are strict. Crossing the $N$th support
exponent means matching its coefficient as well. Merely reaching that
valuation without crossing it need not match the $N$th coefficient.

\begin{corollary}[A degree-limited obstruction]\label{cor:degreebound}
If $a$ is algebraic over $\C(s)$ of degree less than $P_N$, then
\[
 v(h_\delta-a)\le q_{p_N}\delta.
\]
No sequence of algebraic approximants, regardless of the growth of
its degrees, converges to $h_\delta$ in the full valuation topology.
\end{corollary}
\begin{proof}
The first assertion is the contrapositive of
Theorem~\ref{thm:degree}. The second follows from
Theorem~\ref{thm:barrier}: every such approximant stays outside the
valuation neighborhood $v(h_\delta-a)>\delta$.
\end{proof}

\begin{remark}[What the degree law does not measure]
A field degree is not a lower bound for expression length or running
time. The sum $h_{\delta,N}$ has a short expression involving $N$
radicals, despite its primorial degree over $\C(s)$. If the ground field
already contains those radicals, its degree over the new ground field
is one. The sharpness claim concerns exactly the field and precision
specified in Theorem~\ref{thm:degree}.
\end{remark}

\section{Almost-disjoint prime tails}\label{sec:tails}

The single tail suffices for a $2\times2$ counterexample. To hide many
independent negative directions at once, we need a simultaneous
separation estimate that allows arbitrary coefficients from $\Pfield$.
Algebraic independence alone would not give the required valuation bound.

\subsection{An explicit continuum-sized family}

For an infinite set $A\subseteq\primes$, define
\begin{equation}\label{eq:Atail}
 h_A=\sum_{p\in A}t^{q_p\delta}.
\end{equation}
Every such sum is a positive-order real Hahn series. Choose a bijection
from the set of nonempty finite binary words to the primes, listing words
first by length and then lexicographically. For an infinite binary
sequence $b\in\{0,1\}^{\N}$, let $A_b$ consist of the primes assigned to
its nonempty finite prefixes. Each $A_b$ is infinite. Distinct branches
have only finitely many common prefixes, so $A_b\cap A_c$ is finite
when $b\ne c$.

We shall only use the following property of this family: among finitely
many distinct selected sets, each one has infinitely many primes that
belong to none of the other selected sets. The construction makes the
family explicit without choosing a transcendence basis.

\subsection{The fresh-prime separation estimate}

\begin{theorem}[Uniform linear-combination separation]\label{thm:separation}
Let $A_1,\ldots,A_r$ be infinite pairwise almost-disjoint sets of primes,
and write $h_j=h_{A_j}$. If $a_0,a_1,\ldots,a_r\in\Pfield$ and
$(a_1,\ldots,a_r)\ne0$, then
\begin{equation}\label{eq:separation}
 v\!\left(a_0+\sum_{j=1}^r h_j a_j\right)
 <\delta+\min_{1\le j\le r}v(a_j).
\end{equation}
In particular the linear combination on the left is nonzero.
\end{theorem}
\begin{proof}
Let $m_0=\min_{1\le j\le r}v(a_j)$ and choose $j$ with
$v(a_j)=m_0$. Choose one finitely generated subgroup $H$ containing
$\delta$ and every support $\supp a_i$, including $i=0$.
Write $H\cap\Q\delta=(\delta/m)\Z$.

There is a prime $p\in A_j$ that does not divide $m$ and belongs to no
other selected $A_i$. Consider the coefficient at
\[
 \alpha=m_0+q_p\delta.
\]
The term $t^{q_p\delta}a_j$ contributes the nonzero leading coefficient
of $a_j$. No term of $a_0$ contributes, because
$\alpha\notin H$ by Lemma~\ref{lem:cyclicintersection}.
No term $t^{q_q\delta}a_i$ for a different prime $q$ contributes,
because the cosets $q_p\delta+H$ and $q_q\delta+H$ are distinct.
The same prime $p$ occurs in none of the other $h_i$.
Within $t^{q_p\delta}a_j$, only the coefficient at $m_0$ can contribute
to $\alpha$.

Thus the coefficient at $\alpha$ in the entire linear combination is
exactly the nonzero leading coefficient of $a_j$. The Hahn products are
well-defined, so coefficient extraction is legitimate. Consequently the
valuation of the combination is at most
$\alpha=m_0+q_p\delta<m_0+\delta$.
\end{proof}

The argument detects one surviving coefficient rather than estimating
an infinite sum. It does not require a norm, an Archimedean valuation
rank, a cofinal sequence of exponents, or convergence of partial sums.
The exponent $m_0$ can be arbitrarily large, small, or in a different
Archimedean class from $\delta$.

\subsection{Algebraic independence over the entire probe field}

\begin{theorem}[Independent prime-tail family]\label{thm:independence}
For the almost-disjoint family constructed above,
\[
 \{h_{A_b}:b\in\{0,1\}^{\N}\}
\]
is algebraically independent over $\Pfield$.
\end{theorem}
\begin{proof}
We prove algebraic independence of every finite subfamily by induction
on its size. The same argument covers the initial size one.
Suppose $h_1,\ldots,h_{r-1}$ are algebraically independent over
$\Pfield$, and assume that a nonzero polynomial relation involving also
$h_r$ exists. After treating the last variable separately, it becomes
\begin{equation}\label{eq:relation}
 R(h_r)=0,
 \qquad
 R(X)\in\Pfield(h_1,\ldots,h_{r-1})[X],
 \qquad R\ne0.
\end{equation}
Indeed, at least one coefficient of the original polynomial in the last
variable remains nonzero by the induction hypothesis. We may take $R$
with coefficients polynomial, rather than rational, in the previous
$h_j$. Let $d=\deg R\ge1$.

Choose a finitely generated group $H$ containing $\delta$ and all
supports of the finitely many original coefficients from $\Pfield$.
As before, write $H\cap\Q\delta=(\delta/m)\Z$.
Choose a prime
\[
 p\in A_r\setminus\bigcup_{j<r}A_j,
 \qquad p\nmid m,
 \qquad p>d.
\]
Such a prime exists because only finitely many primes are excluded by
the intersections, the divisors of $m$, and the bound $d$.

Let $J$ be the subgroup of $\Gamma$ generated by $H$ and all
$\delta/q$ with $q$ in the selected prime sets and $q\ne p$.
Every rational coefficient of $\delta$ in $J\cap\Q\delta$ has a
reduced denominator not divisible by $p$. To see this, any expression
in $J$ is $h+\ell$, with $h\in H$ and $\ell$ a rational multiple of
$\delta$ whose denominator uses only primes different from $p$.
If the expression lies on the rational line, then so does $h$, and its
denominator divides $m$. Thus $\delta/p\notin J$, whereas
$p(\delta/p)=\delta\in J$. Its class in $\Gamma/J$ has order exactly
$p$.

Choose a primitive $p$th root of unity $\zeta$ and extend the character
that sends $\delta/p+J$ to $\zeta$ to all of $\Gamma/J$.
The associated automorphism $T$ fixes all original coefficients, all
$h_j$ for $j<r$, and every term of $h_r$ except its term at
$q_p\delta=\delta-\delta/p$.
For $k=0,\ldots,p-1$,
\begin{equation}\label{eq:orbit-tail}
 T^k(h_r)=h_r+(\zeta^{-k}-1)t^{q_p\delta}.
\end{equation}
These are $p$ distinct elements. Since $T$ fixes the coefficients of
$R$, they are all roots of the nonzero polynomial $R$. This contradicts
$p>d$. The induction is complete.
\end{proof}

The theorem is stronger than algebraic independence over $\C$ or over a
single rational-function field. Its base contains all monomials and all
series supported in every individual finitely generated subgroup.
The proof uses fresh prime denominators in the \emph{exponents}, rather
than algebraically independent choices of ordinary coefficients.
This distinguishes it from the repository's existing tail-span and
differential-transcendence constructions \cite{repo-map}.

\begin{corollary}\label{cor:trdeg}
For $\Gamma=\Q$,
\[
 \trdeg\bigl(\C((t^\Q))/\mathscr P_\Q\bigr)=2^{\aleph_0}.
\]
\end{corollary}
\begin{proof}
Theorem~\ref{thm:independence} gives the lower bound. A Hahn series on
the countable group $\Q$ is in particular a function $\Q\to\C$.
There are only $|\C|^{\aleph_0}=2^{\aleph_0}$ such functions, giving
the upper bound.
\end{proof}

This cardinal conclusion is included to record what the explicit
construction implies. No novelty is claimed for a generic
transcendence-degree cardinality phenomenon.

\section{Maximal hidden negative inertia}\label{sec:inertia}

We now combine the uniform separation estimate with one finite
congruence calculation. The separation theorem carries the substantive
support information; the matrix algebra is elementary.

\begin{theorem}[Hidden-inertia construction]\label{thm:inertia}
Let $A_1,\ldots,A_r$ be infinite pairwise almost-disjoint prime sets,
$r\ge1$, and put $h_j=h_{A_j}$, $h=(h_1,\ldots,h_r)^T$,
$\eps=t^\delta$. The real symmetric matrix
\[
 A_r=\begin{pmatrix}1&h^T\\h&hh^T-\eps^2I_r\end{pmatrix}
\]
is strictly positive on every nonzero vector in $\Pfield^{r+1}$.
Over $K_\Gamma$, however,
\begin{equation}\label{eq:inertia}
 \In(A_r)=(1,r,0),\qquad \det A_r=(-\eps^2)^r.
\end{equation}
Its standard-part matrix is $\diag(1,0,\ldots,0)$.
\end{theorem}
\begin{proof}
For $z=(a_0,y)$ with $y=(a_1,\ldots,a_r)^T$,
\begin{equation}\label{eq:quadraticidentity}
 z^*A_rz
 =\left|a_0+\sum_{j=1}^r h_ja_j\right|^2
  -\eps^2\sum_{j=1}^r|a_j|^2.
\end{equation}
If $y=0$ and $z\ne0$, this equals $|a_0|^2>0$.
Otherwise set $m_0=\min_jv(a_j)$.
Theorem~\ref{thm:separation} and \eqref{eq:normvaluation} show that the
first term in \eqref{eq:quadraticidentity} is positive and has valuation
strictly less than $2m_0+2\delta$. By \eqref{eq:sum-squares}, the second
term has exactly that valuation. Hence the leading term of the
difference is the positive leading term of the first term.
This proves strict positivity for probe vectors.

For the actual inertia, define
\[
 T=\begin{pmatrix}1&h^T\\0&I_r\end{pmatrix},
 \qquad D=\diag(1,-\eps^2 I_r).
\]
A direct multiplication gives
\begin{equation}\label{eq:congruence}
 A_r=T^*DT,\qquad \det T=1.
\end{equation}
Thus the inertia and determinant are those in \eqref{eq:inertia}.
Equivalently, the $r$-dimensional subspace
\[
 \{(-h^Ty,y):y\in K_\Gamma^r\}
\]
is negative definite, since its quadratic value is
$-\eps^2\sum|y_j|^2$. The vector $e_0$ has value one.
Finally every $h_j$ and $\eps$ has positive valuation, giving the stated
standard part.
\end{proof}

\begin{corollary}[Smallest example]\label{cor:smallmatrix}
With the single all-primes tail $h=h_\delta$,
\begin{equation}\label{eq:smallmatrix}
 A=\begin{pmatrix}1&h\\h&h^2-t^{2\delta}\end{pmatrix}
\end{equation}
is strictly positive on $\Pfield^2\setminus\{0\}$ but has determinant
$-t^{2\delta}$ and value $-t^{2\delta}$ at $(-h,1)$.
\end{corollary}
\begin{proof}
For $b\ne0$, the quadratic value at $(a,b)$ is
\[
 |b|^2\left(|a/b+h|^2-t^{2\delta}\right).
\]
Theorem~\ref{thm:barrier} applies to $-a/b\in\Pfield$.
The remaining statements follow directly from the determinant and
quadratic identity.
\end{proof}

\subsection{Two sharpness statements}

The dimension and inertia counts are optimal in the following simple,
precise senses. In dimension one, testing the vector $1$ detects a
negative scalar, so no example exists. Thus size two is minimal.
In dimension $r+1$, a form that is positive at $e_0$ cannot have
$r+1$ negative directions. Our construction has $r$, the maximum
possible number compatible with even this one positive probe.

These are sharpness claims about dimension and inertia, not about the
ordinal type of the matrix's entire support. Although each $h_j$ has
support of order type $\omega$, products $h_ih_j$ need not have the same
support order type. No support-order threshold for arbitrary matrix
entries is asserted.

\subsection{Why algebraic closedness does not rescue the tests}

The probe field is algebraically closed, but the matrix entries need
not lie in it. Already $h\notin\Pfield$ by
Theorem~\ref{thm:barrier}. A transfer principle for formulas with
parameters from a real-closed subfield does not allow parameters outside
that subfield. There is therefore no conflict with real-closed field
theory.

If the probe field is enlarged to contain the $h_j$, the explicit
negative vectors immediately become available. Similarly, computing the
principal minor on coordinates $0,j$ yields
\[
 \det\begin{pmatrix}1&h_j\\h_j&h_j^2-\eps^2\end{pmatrix}=-\eps^2.
\]
Hence the construction is not an obstruction to an exact positivity
algorithm operating in the field generated by the entries. It is an
obstruction to replacing quantification over that field by a particular
fixed support-restricted collection of test vectors.

\section{Valuation closure and the Levi--Civita case}\label{sec:closure}

Let $\mathscr C_\Gamma$ be the closure of $\Pfield$ inside $K_\Gamma$
for the valuation topology. Explicitly, $a\in\mathscr C_\Gamma$ means
that for every $\rho\in\Gamma$ there exists $b\in\Pfield$ with
$v(a-b)>\rho$. This definition uses neighborhoods, not a presumed
countable cofinal sequence. The closure is a subfield, since addition,
multiplication, and inversion away from zero are continuous in a valued
field.

\begin{theorem}[The barrier persists through valuation closure]
\label{thm:closedseparation}
For the tails in Theorem~\ref{thm:separation}, the estimate
\[
 v\!\left(a_0+\sum_{j=1}^r h_ja_j\right)
 <\delta+\min_{1\le j\le r}v(a_j)
\]
continues to hold for $a_0,\ldots,a_r\in\mathscr C_\Gamma$, provided
$(a_1,\ldots,a_r)\ne0$.
Consequently $A_r$ is strictly positive on every nonzero vector over
$\mathscr C_\Gamma$.
\end{theorem}
\begin{proof}
Suppose the estimate fails and write $m_0=\min_{j\ge1}v(a_j)$.
For every $j=0,\ldots,r$, choose $b_j\in\Pfield$ with
\[
 v(a_j-b_j)>m_0+\delta.
\]
For $j\ge1$, these approximations preserve the minimum valuation:
$\min_jv(b_j)=m_0$. Write
$L_a=a_0+\sum h_ja_j$ and $L_b=b_0+\sum h_jb_j$.
Since each $v(h_j)>0$,
\[
 v(L_a-L_b)>m_0+\delta.
\]
The assumed $v(L_a)\ge m_0+\delta$ therefore implies
$v(L_b)\ge m_0+\delta$, contrary to
Theorem~\ref{thm:separation}. The positivity conclusion follows from
exactly the valuation comparison in Theorem~\ref{thm:inertia}.
\end{proof}

\subsection{Identifying the closure for rational exponents}

\begin{definition}
A subset $S\subseteq\Q$ is \emph{left finite} if for every $\rho\in\Q$
the set $S\cap(-\infty,\rho]$ is finite. The complex Levi--Civita field
is the field of complex Hahn series with left-finite rational support.
\end{definition}

The Puiseux and Levi--Civita relationship belongs to the classical
valuation-theoretic setting; the Puiseux-field conventions are discussed
in \cite{dense-subfields}. We include the elementary closure argument to
make the exact topology explicit.

\begin{proposition}\label{prop:lc}
The closure $\mathscr C_\Q$ of $\mathscr P_\Q$ in
$\C((t^\Q))$ is precisely the complex Levi--Civita field.
\end{proposition}
\begin{proof}
A well-ordered subset of $\frac1n\Z$ is bounded below and has only
finitely many elements below each rational bound. Hence every Puiseux
series has left-finite support.

Suppose $f$ lies in the closure. For each $\rho\in\Q$, take a Puiseux
series $g$ with $v(f-g)>\rho$. The two series have identical
coefficients at all exponents at most $\rho$. The support of $f$ below
that bound is therefore finite. Thus $f$ is left finite.

Conversely, for a left-finite series $f$ and a bound $\rho$, truncate to
the finitely many exponents at most $\rho$. This finite sum is a Puiseux
series, and the difference has valuation greater than $\rho$, or is
zero. Therefore $f$ lies in the closure.
\end{proof}

\begin{corollary}\label{cor:lc}
There are explicit Hermitian matrices over $\C((t^\Q))$ with inertia
$(1,r,0)$ that are strictly positive on every nonzero vector with
complex Levi--Civita coordinates.
\end{corollary}

The coefficients of these matrices are \emph{not} all in the
Levi--Civita field. In particular, $h_\delta$ has infinitely many
exponents below $\delta$ when $\delta\in\Q_{>0}$.
The corollary does not say that a Hermitian matrix defined over the
Levi--Civita field can be positive there and become indefinite under an
ordered extension. It says that Levi--Civita probes are insufficient for
certain matrices defined in a larger Hahn field.

\section{Density and actual surreal numbers}\label{sec:density}

\subsection{The elementary density principle}

The following proposition explains the general mechanism independently
of generalized series. It is included as an elementary organizing fact,
not as a new theorem of ordered-field theory.

\begin{proposition}[Quadratic testing characterizes order density]
\label{prop:density}
Let $E\subseteq F$ be ordered fields. The following conditions are
equivalent.
\begin{enumerate}[label=(\roman*)]
\item $E$ is order dense in $F$.
\item For every finite symmetric matrix over $F$, nonnegativity on all
vectors in $E^n$ implies nonnegativity on all vectors in $F^n$.
\item The implication in (ii) holds for matrices of size two.
\end{enumerate}
If $E$ is not order dense, a size-two form can in fact be strictly
positive on every nonzero $E$-vector and negative at an $F$-vector.
\end{proposition}
\begin{proof}
Assume density. If a quadratic form is negative at a point in $F^n$,
polynomial continuity in the product order topology makes it negative
on some product of open intervals around that point. Each interval
contains an element of $E$. Thus an $E$-vector already has negative
value. This proves (i)$\Rightarrow$(ii), and
(ii)$\Rightarrow$(iii) is immediate.

Conversely, if $E$ is not dense, choose $a<b$ in $F$ such that
$(a,b)\cap E=\varnothing$. Put
\[
 h=\frac{a+b}{2},\qquad e=\frac{b-a}{4}>0.
\]
The closed interval $[h-e,h+e]$ lies inside $(a,b)$.
For
\[
 q(x,y)=(x-hy)^2-e^2y^2,
\]
if $(x,y)\in E^2$ is nonzero and $y=0$, then $q=x^2>0$.
If $y\ne0$, the ratio $x/y\in E$ is outside that closed interval,
so $|x/y-h|>e$ and again $q>0$. Yet $q(h,1)=-e^2<0$.
This disproves (iii).
\end{proof}

For the corresponding complexified fields $E[i]\subseteq F[i]$, the
same conclusion follows by applying continuity to real and imaginary
coordinates. The real witness above also works directly: if
$z/w=x+iy$ with $x,y\in E$, then
$|z/w-h|^2=(x-h)^2+y^2>e^2$.

This proposition is why a bare example over a nondense subfield would
be too weak a contribution by itself. The prime-tail package specifies
a particularly large natural probe field, gives a precise missed
valuation scale, defeats its closure, constructs the maximal number of
negative directions, and computes exact algebraic approximation costs.
General results about dense subfields of henselian fields address a
broader but different question \cite{dense-subfields}; nothing here says
that a Hahn field has no proper dense subfields.

\subsection{Transport to Conway normal forms}

If $\Gamma$ is a set-sized ordered additive subgroup of $\No$, the
normal-form correspondence
\begin{equation}\label{eq:transport}
 \sum_{\gamma\in S}c_\gamma t^\gamma
 \longmapsto
 \sum_{\gamma\in S}c_\gamma\omega^{-\gamma}
\end{equation}
embeds its real Hahn field into $\No$ and its complex Hahn field into
$\No[i]$. The minus sign is essential: increasing Hahn exponents become
decreasing Conway growth exponents. The correspondence preserves
addition, multiplication, conjugation, and the leading-coefficient
order. These are the standard normal-form operations
\cite{gonshor}; the repository also emphasizes this orientation issue
\cite{repo-map}.

Take $\Gamma=\Q$ and $\delta=1$. Then
\begin{equation}\label{eq:surrealtail}
 h=\sum_{p\in\primes}\omega^{-1+1/p}
 =\omega^{-1/2}+\omega^{-2/3}+\omega^{-4/5}
  +\omega^{-6/7}+\cdots
\end{equation}
is an actual surreal number. No limiting interpretation is required;
this is its Conway normal form. The actual surcomplex matrix
\begin{equation}\label{eq:surrealmatrix}
 \begin{pmatrix}1&h\\h&h^2-\omega^{-2}\end{pmatrix}
\end{equation}
has a negative quadratic value at $(-h,1)$ and is strictly positive on
all nonzero finite-lattice-supported surcomplex probe vectors.

\begin{theorem}[Workspace-independent surreal probe obstruction]
\label{thm:surreal}
The positivity conclusion for \eqref{eq:surrealmatrix}, and the
higher-dimensional construction transported by \eqref{eq:transport},
holds for every finite surcomplex vector whose coordinates each have
Conway support contained in some finitely generated additive subgroup
of $\No$. The subgroup may depend on the vector and may contain
arbitrarily large surreal exponents.
\end{theorem}
\begin{proof}
For the given finite vector, choose a finitely generated subgroup $H$
containing all its exponent supports and $1$. Its divisible hull is a
set-sized ordered abelian group $\Gamma'\subseteq\No$.
Both the vector and the prime-tail matrix are represented in the Hahn
workspace for $\Gamma'$, with Hahn exponents opposite to the Conway
exponents. The probe coordinates lie in $\mathscr P_{\Gamma'}$.
Theorem~\ref{thm:inertia} applies there and transports back.
\end{proof}

Thus adding any finitely many new monomial scales to a particular probe
representation does not remove the obstruction. The quantification is
local to each finite input. No set containing all surreal numbers, and
no class-indexed Hahn summation, is being used.

An analogous statement applies to actual surcomplex coordinates
algebraic over rational expressions in surreal monomials. Any particular
algebraic equation uses finitely many such expressions, hence finitely
many monomial generators. The character support bound places its roots
in a finite-lattice workspace. Again this does not include a coordinate
whose defining parameters already contain the prime tail itself.

\section{The exact two-scale matrix criterion}\label{sec:twoscale}

We now turn to the second part of the paper. Here $\Gamma$ may be any
nonzero ordered abelian group; divisibility is unnecessary. Let
$\eps=t^\eta$ for some $\eta>0$.
The coefficients of the two matrices below are ordinary complex
numbers, while vectors in the positivity test are allowed to have
arbitrary Hahn coordinates.

\subsection{An isotropic vector of a positive form lies in its kernel}

\begin{lemma}\label{lem:psdkernel}
Let $L$ be an ordered field, let $A=A^*$ over $L[i]$, and assume
$x^*Ax\ge0$ for all $x$. If $u^*Au=0$, then $Au=0$.
\end{lemma}
\begin{proof}
Fix a vector $w$ and write $b=w^*Au$ and $c=w^*Aw\ge0$.
If $b\ne0$, choose $s=-\lambda b$ with $\lambda>0$ in $L$.
Then
\[
 (u+sw)^*A(u+sw)
 =-2\lambda|b|^2+\lambda^2|b|^2c.
\]
For $c=0$ this is negative. For $c>0$, taking $\lambda=1/c$
also makes it negative. Both contradict positivity. Hence
$w^*Au=0$ for every $w$, and the coordinate vectors give $Au=0$.
\end{proof}

\begin{theorem}[Two-scale matrix positivity]\label{thm:twoscale}
Let $P,Q\in M_d(\C)$ be Hermitian. If $P\succeq0$, let
$N=\ker P$ and let $C_N$ denote the compression of $Q$ to $N$,
viewed as an ordinary Hermitian operator on $N$.
Then
\begin{equation}\label{eq:twoscalecriterion}
 P+\eps Q\succeq0
 \quad\Longleftrightarrow\quad
 \begin{cases}
 P\succeq0,\\
 C_N\succeq0,\\
 Q u=0\quad\text{for every }u\in\ker C_N\subseteq N.
 \end{cases}
\end{equation}
Equivalently, the last two clauses say that the quadratic form of $Q$
is nonnegative on $N$, and that
\begin{equation}\label{eq:annihilationcondition}
 u\in N,\quad u^*Qu=0\quad\Longrightarrow\quad Qu=0.
\end{equation}
When these conditions hold,
\begin{align}
 \ker(P+\eps Q)&=(\ker P\cap\ker Q)\otimes_\C K_\Gamma,
 \label{eq:kernel}\\
 \rank(P+\eps Q)&=\rank P+\rank C_N.
 \label{eq:rank}
\end{align}
\end{theorem}
\begin{proof}
Suppose $P+\eps Q\succeq0$. Testing an ordinary vector whose $P$-value
is negative would give a negative leading coefficient, so $P\succeq0$.
For ordinary $u\in N$, the value is $\eps u^*Qu$, proving
$C_N\succeq0$. If $u\in\ker C_N$, then $u^*Qu=0$.
Its value under $P+\eps Q$ is zero, so
Lemma~\ref{lem:psdkernel} gives $(P+\eps Q)u=0$.
As $Pu=0$, this means $Qu=0$.

Conversely assume the three clauses. Use an ordinary unitary change of
basis to decompose
\[
 \C^d=R\oplus N_+\oplus N_0,
 \qquad R=\operatorname{ran}P,
 \quad N_+=\operatorname{ran}C_N,
 \quad N_0=\ker C_N.
\]
The matrices $P_R$ and $C_+=\restr{C_N}{N_+}$ are positive definite on
their respective spaces. By hypothesis $Q$ annihilates $N_0$, so the
entire matrix is zero on that last direct summand. On $R\oplus N_+$ it
has the block form
\begin{equation}\label{eq:activeblock}
 \begin{pmatrix}
 P_R+\eps A&\eps B\\
 \eps B^*&\eps C_+
 \end{pmatrix}.
\end{equation}
Let $U=P_R+\eps A$. Every leading principal minor of $U$ has positive
constant coefficient, so $U$ is positive definite over the ordered
Hahn field. This is also immediate from an $LDL^*$ decomposition.
Its inverse has entries rational in $\eps$ and constant term $P_R^{-1}$.
The Schur complement of $U$ in \eqref{eq:activeblock} is
\[
 \eps\bigl(C_+-\eps B^*U^{-1}B\bigr).
\]
Every leading principal minor of the factor in parentheses has the
positive constant coefficient of the corresponding minor of $C_+$.
That factor, and hence the Schur complement, is positive definite.
Block Gaussian congruence now proves positivity on the entire active
space. If either active space has dimension zero, the same argument
reduces to its remaining positive block without taking an empty inverse.

The only kernel is therefore $N_0\otimes_\C K_\Gamma$.
Within $N$, the condition $Qu=0$ is equivalent to $u\in\ker C_N$
under our hypotheses. This proves \eqref{eq:kernel} and
\eqref{eq:rank}. All inverses used are rational expressions in $\eps$;
no sequential convergence assertion is involved.
\end{proof}

The theorem is a finite Schur-complement criterion. Its role is to
expose the exact null-direction coupling needed by the measure theorem.
We do not claim Schur complements or singular perturbation theory as
new techniques.

\subsection{The two ways a residual direction behaves}

\begin{example}[A coupling that positivity forbids]\label{ex:badcoupling}
Take
\[
 P=\begin{pmatrix}1&0\\0&0\end{pmatrix},\qquad
 Q=\begin{pmatrix}0&1\\1&0\end{pmatrix}.
\]
Here $N=\C e_2$ and $C_N=0$, but $Qe_2=e_1\ne0$.
Thus
\[
 P+\eps Q=\begin{pmatrix}1&\eps\\\eps&0\end{pmatrix}
\]
is not positive: its determinant is $-\eps^2$, and $(-\eps,1)$ has
value $-\eps^2$.
Nevertheless every ordinary vector $u=(a,b)$ has nonnegative value:
if $a\ne0$, its leading value is $|a|^2>0$, and if $a=0$, its value is
zero. Ordinary scalar vector tests miss the coupling completely.
\end{example}

\begin{example}[A coupling that positivity allows]\label{ex:goodcoupling}
With the same $P$ but
\[
 Q=\begin{pmatrix}0&1\\1&1\end{pmatrix},
\]
the compression $C_N$ is strictly positive. Then
\[
 \begin{pmatrix}1&\eps\\\eps&\eps\end{pmatrix}\succ0,
 \qquad \det=\eps-\eps^2>0.
\]
The ordinary matrix $Q$ itself is indefinite. Requiring $Q\succeq0$
would therefore be too strong. Only its action on the leading kernel
and the coupling out of its residual zero space are relevant.
\end{example}

\begin{example}[Why more than two scales changes the question]
\label{ex:thirdscale}
For ordinary real $c$,
\[
 B_c=\begin{pmatrix}1&\eps\\\eps&c\eps^2\end{pmatrix}
\]
is congruent to $\diag(1,(c-1)\eps^2)$ and hence is positive
semidefinite exactly when $c\ge1$. For every $c\ge0$, however, all
ordinary vector tests are nonnegative: a vector with first coordinate
nonzero has positive leading value, and one with first coordinate zero
has value $c\eps^2|b|^2$.

The third-scale coefficient can repair, exactly cancel, or fail to repair
the second-order Schur correction. Therefore Theorem~\ref{thm:twoscale}
is not an arbitrary-support matrix criterion in disguise. Its
annihilation condition is forced because no third-scale coefficient is
available to repair the coupling.
\end{example}

\section{A two-scale matrix-valued null-ideal theorem}\label{sec:measure}

\subsection{Definitions and the quantifiers}

Let $(X,\Sigma)$ be a measurable space. A finite Hermitian
$\C^{d\times d}$-valued measure is a matrix of finite complex measures
$M=(M_{ij})$ with $M(E)^*=M(E)$ for every measurable $E$.
Countable additivity is entrywise, and each entry has finite total
variation. No topology on $X$ or regularity assumption is needed.

For two such measures, define
\[
 M=M_0+\eps M_1.
\]
This is a coefficientwise Hahn-valued matrix measure with two possible
exponents. We call it \emph{positive} when
\begin{equation}\label{eq:measurepositive}
 z^*M(E)z\ge0
 \quad\text{for every }E\in\Sigma
 \text{ and every }z\in K_\Gamma^d.
\end{equation}
The vector $z$ may depend on $E$ and may contain infinitesimals.
The model is coefficientwise countable additivity. Strong Hahn
summability of all countable families of event masses is not assumed.
The repository distinguishes these notions, and we retain that
separation \cite{repo-herglotz,repo-map}.

If $M_0$ is an ordinary positive matrix measure, then for
$u\in\C^d$
\begin{equation}\label{eq:scalarnotations}
 \mu_u(E)=u^*M_0(E)u
\end{equation}
is a finite positive real measure. For $M_1$, set
\begin{equation}\label{eq:tailnotations}
 \nu_u(E)=u^*M_1(E)u,
 \qquad \sigma_u(E)=M_1(E)u.
\end{equation}
The first is a finite signed real measure and the second a finite
complex vector measure. Write $\nu_u=\nu_u^+-\nu_u^-$ for its Jordan
decomposition, and $|\sigma_u|$ for total variation using the ordinary
Euclidean norm. Any other ordinary norm on $\C^d$ gives the same null
sets.

For a positive scalar measure $\lambda$, the assertion
$|\sigma_u|\ll\lambda$ means that $\sigma_u$ vanishes on every
measurable subset of every $\lambda$-null set. Indeed, this is equivalent
to the usual variation formulation by the definition of total
variation through finite partitions.

\subsection{The characterization}

\begin{theorem}[Two-scale matrix null ideals]\label{thm:measure}
Let $M_0,M_1$ be finite Hermitian matrix-valued measures and
$\eps=t^\eta$ with $\eta>0$. Then $M_0+\eps M_1$ is positive in the
sense of \eqref{eq:measurepositive} if and only if all of the following
hold:
\begin{enumerate}[label=(\roman*)]
\item $M_0$ is an ordinary positive matrix-valued measure.
\item For every ordinary $u\in\C^d$,
\begin{equation}\label{eq:scalar-null}
 \nu_u^-\ll\mu_u.
\end{equation}
\item For every ordinary $u\in\C^d$,
\begin{equation}\label{eq:vector-null}
 |\sigma_u|\ll\mu_u+\nu_u^+.
\end{equation}
\end{enumerate}
The result holds for every nonzero ordered value group $\Gamma$;
divisibility is not required. No bounded Radon--Nikodym derivative is
required in either absolute-continuity condition.
\end{theorem}
\begin{proof}
\emph{Necessity.}
Assume $M$ is positive. For every ordinary $u$ and event $E$, a negative
value $u^*M_0(E)u$ would be the negative leading coefficient of
$u^*M(E)u$. Thus $M_0$ is positive.

Fix $u$ and let $E$ be $\mu_u$-null. Every measurable $B\subseteq E$
is also $\mu_u$-null. Positivity gives
\[
 0\le u^*M(B)u=\eps\nu_u(B),
\]
so $\nu_u$ restricts to a positive measure on $E$.
Therefore $\nu_u^-(E)=0$, proving \eqref{eq:scalar-null}.

Next suppose $(\mu_u+\nu_u^+)(E)=0$. By
\eqref{eq:scalar-null}, $\nu_u^-(E)=0$ as well.
For every measurable $B\subseteq E$, we consequently have
\[
 u^*M_0(B)u=0,\qquad u^*M_1(B)u=0.
\]
The ordinary matrix $M_0(B)$ is positive, so
$M_0(B)u=0$ by Lemma~\ref{lem:psdkernel}. The Hahn matrix $M(B)$ is
positive and has zero quadratic value at $u$, so $M(B)u=0$ too.
Subtracting yields $\eps M_1(B)u=0$, whence $\sigma_u(B)=0$.
Thus the vector measure vanishes on every subset of $E$, which means
$|\sigma_u|(E)=0$. This proves \eqref{eq:vector-null}.

\emph{Sufficiency.}
Assume (i)--(iii), and fix an event $E$. Put $P=M_0(E)$ and
$Q=M_1(E)$. The matrix $P$ is positive. If an ordinary vector
$u\in\ker P$, then $\mu_u(E)=0$, and (ii) gives
$\nu_u^-(E)=0$. Hence
\[
 u^*Qu=\nu_u(E)=\nu_u^+(E)\ge0.
\]
The compression of $Q$ to $\ker P$ is therefore positive.
If in addition $u^*Qu=0$, then $\nu_u^+(E)=0$.
The measure $\mu_u+\nu_u^+$ vanishes on $E$, so (iii) gives
$|\sigma_u|(E)=0$ and in particular $Qu=\sigma_u(E)=0$.
The exact hypotheses of Theorem~\ref{thm:twoscale} are now satisfied.
It follows that $M(E)=P+\eps Q$ is positive on every Hahn vector.
Since $E$ was arbitrary, $M$ is positive.
\end{proof}

The theorem is a null-ideal characterization, not a finite test in $u$.
It quantifies over all ordinary vectors and all measurable null sets.
Its content is that these ordinary measure-theoretic conditions are
exactly sufficient to handle the vastly larger family of Hahn-valued
vectors in \eqref{eq:measurepositive}.

\begin{corollary}[Eventwise kernel and rank]\label{cor:measurekernel}
Under the equivalent conditions of Theorem~\ref{thm:measure}, for every
event $E$,
\[
 \ker M(E)=\bigl(\ker M_0(E)\cap\ker M_1(E)\bigr)
                  \otimes_\C K_\Gamma.
\]
If $C_E$ is the compression of $M_1(E)$ to $\ker M_0(E)$, then
\[
 \rank M(E)=\rank M_0(E)+\rank C_E.
\]
\end{corollary}
\begin{proof}
Apply the kernel and rank clauses of Theorem~\ref{thm:twoscale} to
each event separately.
\end{proof}

This corollary is eventwise. It is not a claim of a common kernel for all
events, nor of a spectral measure theorem for an infinite-dimensional
Hahn Hilbert space.

\subsection{Reduction to the known scalar criterion}

\begin{corollary}[Scalar specialization]\label{cor:scalarmeasure}
For finite real signed measures $\mu_0,\mu_1$, the coefficientwise
Hahn measure $\mu_0+\eps\mu_1$ is positive if and only if
\[
 \mu_0\ge0,\qquad \mu_1^-\ll\mu_0.
\]
In dimension one, condition \eqref{eq:vector-null} adds no restriction.
\end{corollary}
\begin{proof}
For $u\ne0$, $\mu_u=|u|^2\mu_0$,
$\nu_u^\pm=|u|^2\mu_1^\pm$, and
$|\sigma_u|=|u|\,|\mu_1|$. If $\mu_0+\mu_1^+$ vanishes on an event,
then $\mu_1^-$ also vanishes there by its absolute continuity with
respect to $\mu_0$. Hence $|\mu_1|$ vanishes there, proving (iii).
The case $u=0$ is trivial.
\end{proof}

This specialization is the scalar result already in the repository
\cite{repo-herglotz}. The new matrix ingredient is not scalarization
itself, but the extra vector-variation condition that prevents a
residual zero direction from coupling to a positive leading direction.

\subsection{Examples showing the role of each condition}

\begin{example}[All scalar null tests pass, but matrix positivity fails]
\label{ex:measurebad}
Let $\rho$ be any nonzero finite positive measure and use the matrices
$P,Q$ from Example~\ref{ex:badcoupling}. Set $M_0=P\rho$ and
$M_1=Q\rho$. For $u=(a,b)$,
\[
 \mu_u=|a|^2\rho,\qquad
 \nu_u=2\operatorname{Re}(\bar a b)\rho.
\]
If $a\ne0$, then $\nu_u^-\ll\mu_u$; if $a=0$, then $\nu_u=0$.
Thus condition (ii) holds for every $u$.
For $u=e_2$, however,
\[
 \mu_u=\nu_u^+=0,
 \qquad \sigma_u=e_1\rho\ne0.
\]
Condition (iii) fails. On any event of positive $\rho$-measure, the
Hahn vector $(-\eps,1)$ detects a negative value.
This works even on a one-point space, so the missing condition is not a
subtlety of infinite measure spaces.
\end{example}

\begin{example}[An indefinite coefficient is harmless]
Use instead $Q$ from Example~\ref{ex:goodcoupling} and again put
$M_0=P\rho$, $M_1=Q\rho$. The measure is positive, because each event
matrix is a nonnegative ordinary multiple of a positive definite
Hahn matrix. For $u=(a,b)$ with $a\ne0$, the leading control measure is
already a positive multiple of $\rho$. For $a=0$ and $b\ne0$,
$\nu_u^+=|b|^2\rho$ supplies the control that (iii) needs.
Thus the vector condition permits the coupling, even though $Q$ is
indefinite as an ordinary matrix.
\end{example}

\begin{example}[No bounded density is necessary]\label{ex:unboundeddensity}
On $(0,1]$ with Lebesgue measure, take
\[
 M_0=I_d\,dx,\qquad
 M_1=-x^{-1/2}I_d\,dx.
\]
Every entry has finite total variation. The density $x^{-1/2}$ is
unbounded, but for each nonzero ordinary $u$, both the scalar negative
part and the vector variation are absolutely continuous with respect
to $\mu_u=\|u\|^2dx$.
Thus $M_0+\eps M_1$ is positive.
Directly, if an event has positive measure, its matrix mass is
\[
 \left(|E|-\eps\int_E x^{-1/2}\,dx\right)I_d\succ0;
\]
if it has zero measure, both coefficients vanish. The infinitesimal
comparison is made separately for each event; no ordinary uniform bound
on density ratios is needed.
\end{example}

\section{What has been established, and what remains}\label{sec:scope}

\subsection{Two distinct hierarchies of probes}

The simplest coupling example
$\begin{psmallmatrix}1&\eps\\\eps&0\end{psmallmatrix}$ has coefficients
in a one-generator Hahn subfield. Ordinary complex vectors miss its
negativity, but a finite-lattice vector detects it immediately.
The prime-tail matrix is more severe: every finite-lattice vector, every
vector algebraic over all monomials, and every vector in the valuation
closure of the finite-lattice field misses it. A vector using the tail
itself detects it immediately.

These are different levels of restriction. They should not be collapsed
into a statement that ``finite calculations cannot detect positivity.''
In both examples an exact $2\times2$ determinant does detect the failure.
What changes is the support needed in a \emph{quadratic witness}, not the
number of rows or the degree of a determinant.

\subsection{Results and imported inputs}

The proofs give a short dependency structure. Character extension gives
the algebraic support-index bound. The support-index bound yields the
exact primorial degree law and the algebraic closedness of the probe
field. A fresh-prime coset argument gives the stronger simultaneous
valuation separation. One congruence and a leading-term comparison then
give maximal hidden negative inertia. The closure theorem is an
approximation contradiction and does not require a new support lemma.

The matrix-measure theorem is logically independent of the prime-tail
construction. Its only algebraic ingredient is the exact two-scale
Schur-complement criterion, and its measure-theoretic ingredient is the
ordinary characterization of a vector measure's null sets by its total
variation. The two parts meet conceptually in the difference between
ordinary vector tests and all Hahn vector tests.

\subsection{Questions genuinely left by this work}

\begin{question}[More than two measure scales]
For a finite or transfinite well-ordered family of ordinary Hermitian
matrix measures $\sum_\gamma t^\gamma M_\gamma$, is there an equally
intrinsic null-ideal characterization that incorporates the Schur
corrections generated by off-diagonal coefficients?
\end{question}

Theorem~\ref{thm:measure} answers this only when there are two scales.
Example~\ref{ex:thirdscale} explains why simply repeating its
annihilation condition at each coefficient cannot be correct: later
coefficients can compensate for earlier off-diagonal effects.
No answer for arbitrary supports is claimed here.

\begin{question}[Optimal probe classes under prescribed support closure]
Given a support-family subfield $E\subseteq F_\Gamma$, which closure
conditions on its supports imply order density, and which quantitative
support invariants bound the degree of approximants to an omitted
valuation ball?
\end{question}

Proposition~\ref{prop:density} reduces qualitative quadratic testing to
density. The prime-tail construction supplies one exact quantitative
answer. It does not classify all support families.

\begin{question}[Relative precision costs]
How does the exact approximation law change when the ground field
already contains an infinite prescribed collection of monomial roots,
or when an algebraic approximant is charged by a tower of extensions
rather than its total field degree?
\end{question}

The present theorem deliberately fixes the ground field $\C(t^\delta)$.
A general extension-degree complexity theory would require a separate
cost model and should not be inferred from the primorial formula.

\subsection{Non-claims}

No result here gives a new construction of the surreal field, a new
spectral theorem over real-closed fields, or a general solution to a
published surreal conjecture. The algebraic independence theorem is
ordinary algebraic independence, not differential independence for the
Berarducci--Mantova derivation. No integration on the full proper-class
surcomplex plane, no positive spectral measure representation, and no
infinite-dimensional Hilbert completion are constructed.

The valuation topology used for closures is the topology of a fixed
set-sized Hahn workspace. It is not the full fine topology on all of
$\No$ or $\No[i]$. The actual-surreal corollaries use normal-form
transport, not a claim that countable partial sums converge in that
fine topology.

Finally, the novelty assessment is a research judgment with limited
search coverage. The written proofs can be checked independently of
that assessment. A future discovery of a prior formulation would change
the priority claim without changing the validity or usefulness of the
explicit constructions.

\appendix
\section{Repository and literature audit}\label{app:audit}

\subsection{Pinned repository scope}

The repository comparison uses
\begin{center}
\small\pin
\end{center}
for \url{https://github.com/VladimirReshetnikov/Surreal}.
The inspection was performed on 22 September 2026. The repository may
have changed after that snapshot.

The reader map at this pin describes 36 reports and stresses that the
collection consists of research drafts, with separate mathematical,
finite-computational, and Lean-verification evidence. We inspected that
map and the relevant spectral and Herglotz report guides, together with
the dynamics guide to exclude another initially considered direction.
We did not read every manuscript and every Lean file line by line.

\begin{longtable}{@{}p{.29\textwidth}p{.64\textwidth}@{}}
\toprule
Inspected item & Consequence for this manuscript\\
\midrule
\endhead
\path{docs/README.md}
& Existing topics include finite and infinite spectral theory, Hahn
measures, Herglotz positivity, tail-span differential transcendence,
normal forms, and several analytic theories. These subjects are not
presented here as missing wholesale.\\[5pt]
\path{spectral-theory/README.md}
& Hermitian diagonalization, finite congruence, Schur corrections,
finite-index monomial extensions, and character calculations already
occur. The present manuscript uses these ideas but claims neither their
invention nor their absence from the repository.\\[5pt]
\path{hahn-herglotz-positivity/README.md}
& Contains the scalar null-ideal criterion, two-scale scalar
absolute-continuity condition, matrix halo normalization, and
negative-mass examples. Explicitly excludes a matrix-valued null-ideal
criterion. Theorem~\ref{thm:measure} supplies the two-scale matrix case,
not an arbitrary-support replacement.\\[5pt]
\path{dynamics-and-normal-forms/README.md}
& Already contains extensive multiplier and surreal-frequency
small-divisor results. No new result in those directions is claimed
here.\\[5pt]
Targeted code search
& Searches for ``algebraic probes'' and ``finite-lattice'' returned no
matches. Such negative searches are limited by terminology and indexing;
they do not establish absence of an equivalent theorem.\\
\bottomrule
\end{longtable}

The shortened report paths in the table are under
\texttt{docs/surcomplex/}. The long report guides were not uniformly
returned in full by the connector; the conclusions above are confined
to the returned relevant portions. In particular, this was a targeted
scope audit, not a repository-wide theorem-equivalence check.

\subsection{Literature positioning}

Kuhlmann and Serra's work on Hahn automorphisms supplies the established
context for diagonal character actions \cite{automorphisms}.
Krapp, Kuhlmann, and Serra study support families defining subfields
\cite{rayner-structures}, in the tradition of Rayner
\cite{rayner}. Kuhlmann's work on dense subfields of henselian fields
provides broader density context \cite{dense-subfields}.
Gonshor supplies the classical surreal normal-form setting
\cite{gonshor}.

The automorphism paper was inspected in its arXiv PDF; the support-family
and dense-subfield papers were inspected in available arXiv text.
The original Rayner paper and the full Gonshor book were not audited
line by line; they are cited for established background. Additional
keyword searches did not locate the exact package of statements here,
but some searches returned low-relevance results. Such searches are not
evidence strong enough to certify priority.

The candidates for independent originality are the exact primorial
precision theorem, the almost-disjoint-tail separation over the full
finite-lattice field and its closure, their maximal-hidden-inertia
application, and the vector-total-variation formulation of the
matrix-measure criterion. The elementary density equivalence and the
support-field construction are intentionally outside that claim.

\section{Finite verification and proof boundaries}\label{app:checks}

The archive contains \texttt{code/verify.py}. It uses exact integer,
rational, and symbolic calculations; it does not substitute a small
floating-point number for an infinitesimal. Its captured output is in
\texttt{data/verification.json} and \texttt{data/verification.txt}.
The recorded run passed all tests: 5946 matrix pairs, of which 1823 were
positive, and 35100 exact fresh-prime nonmembership comparisons.
The script checks the following finite statements.

\begin{enumerate}
\item For $1\le r\le4$, the congruence identity
$A_r=T^*\diag(1,-\eps^2I_r)T$, determinant, and explicit negative
vectors, using independent symbolic real $h_j$.
\item For $1\le N\le6$, the cyclic character stabilizer of the first
$N$ prime-tail exponents modulo $P_N$. It is trivial in each case,
so the finite orbit has size $P_N$.
\item Fresh-prime coset separation for bounded denominators
$1\le m\le60$ and primes at most $97$, with exact rational arithmetic.
\item Every pair consisting of a diagonal $0$--$1$ matrix $P$ and a real
symmetric matrix $Q$ with entries in $\{-1,0,1\}$ in dimensions
$1,2,3$. There are $6+108+5832=5946$ pairs. It compares
Theorem~\ref{thm:twoscale} with the signs of all principal minors of
$P+tQ$, using their first nonzero polynomial coefficient.
\item The determinant identities in the forbidden-coupling,
allowed-coupling, and three-scale repair examples.
\end{enumerate}

The exhaustive test in item 4 covers only those finite real matrices.
It does not establish the theorem for arbitrary complex matrices,
arbitrary dimension, or arbitrary Hahn value groups. Item 2 checks
finite stabilizers, not the infinite algebraic-independence theorem.
No finite test here proves a claim about all measurable sets, all
ordinary vectors, all Hahn supports, or all surreal workspaces.
Those quantifiers are handled by the written proofs.

\subsection{A possible formalization order}

A future Lean development could begin with the finite congruence and
the two-scale algebra lemma, continue with the scalar/vector null-set
measure equivalence, and then develop support-index and fresh-prime
lemmas. The latter require a precise API for character extension,
coefficientwise Hahn automorphisms, and finite generation on a rational
line. Existing repository support, complexification, and spectral
components should be reused rather than duplicated.

This is an implementation outline only. No Lean source or claim of
machine-checked coverage for the new theorems accompanies this archive.

\begin{thebibliography}{99}
\small\raggedright

\bibitem{gonshor}
H.~Gonshor,
\emph{An Introduction to the Theory of Surreal Numbers},
London Mathematical Society Lecture Note Series 110,
Cambridge University Press, 1986.
\href{https://www.cambridge.org/core/books/an-introduction-to-the-theory-of-surreal-numbers/312AE504A3E88E804054BFB390446374}{Publisher record}.

\bibitem{neumann}
B.~H.~Neumann,
On ordered division rings,
\emph{Transactions of the American Mathematical Society}
\textbf{66} (1949), 202--252.
\url{https://doi.org/10.1090/S0002-9947-1949-0032593-5}.

\bibitem{rayner}
F.~J.~Rayner,
An algebraically closed field,
\emph{Glasgow Mathematical Journal} \textbf{9} (1968), 146--151.
\url{https://doi.org/10.1017/S0017089500000422}.

\bibitem{rayner-structures}
L.~S.~Krapp, S.~Kuhlmann, and M.~Serra,
On Rayner structures,
\emph{Communications in Algebra} \textbf{50} (2022), no.~3, 940--948.
\url{https://doi.org/10.1080/00927872.2021.1976789}.
Preprint: \url{https://arxiv.org/abs/2004.03239}.

\bibitem{automorphisms}
S.~Kuhlmann and M.~Serra,
The automorphism group of a valued field of generalised formal power
series,
\emph{Journal of Algebra} \textbf{605} (2022), 339--376.
\url{https://doi.org/10.1016/j.jalgebra.2022.04.023}.
Preprint, version 3: \url{https://arxiv.org/abs/2107.03362v3}.

\bibitem{dense-subfields}
F.-V.~Kuhlmann,
Dense subfields of henselian fields, and integer parts,
in \emph{Logic in Tehran}, Lecture Notes in Logic 26,
Association for Symbolic Logic, 2006, 204--226.
Preprint: \url{https://arxiv.org/abs/1003.5681}.

\bibitem{repo-map}
V.~Reshetnikov, \repo\ repository,
\emph{The surreal and surcomplex reports},
\path{docs/README.md}, snapshot \pin,
accessed 22 September 2026.
\href{https://github.com/VladimirReshetnikov/Surreal/blob/048b72cf7cbfc8ab246e4f73788c10460cb3f6e0/docs/README.md}{Pinned repository reader map}.

\bibitem{repo-spectral}
\repo\ research collection,
\emph{Surcomplex Spectral Theory: Singular Values, Infinitesimal Rank,
Multiscale Stability, Exact Determinantal Ramification, and Hankel
Square-Class Profiles}, report guide, snapshot \pin.
\href{https://github.com/VladimirReshetnikov/Surreal/blob/048b72cf7cbfc8ab246e4f73788c10460cb3f6e0/docs/surcomplex/spectral-theory/README.md}{Pinned spectral-theory report guide}.

\bibitem{repo-herglotz}
\repo\ research collection,
\emph{Positivity without a Positive Measure: Hahn--Herglotz
Normalization and Strict Moment Hierarchies in Surcomplex Analysis},
report guide, snapshot \pin.
\href{https://github.com/VladimirReshetnikov/Surreal/blob/048b72cf7cbfc8ab246e4f73788c10460cb3f6e0/docs/surcomplex/hahn-herglotz-positivity/README.md}{Pinned Hahn--Herglotz report guide}.

\end{thebibliography}
\end{document}
